Siempre me han parecido muy interesantes los avances tecnológicos, y también los riesgos a los que nos exponemos con ellos sin siquiera saberlo. Cuando descubrí que a través de un simple enlace es posible inyectar código o extraer datos de un dispositivo, me pareció algo sorprendente. Recuerdo haber pensado en lo ingenioso que tenía que ser alguien para crear algo así, y en lo extraño que era que todo ese ingenio se usara para hacer daño.

Desde entonces adopté una costumbre: cada vez que no estoy segura de un enlace, lo paso por VirusTotal antes de abrirlo. Un día, en una conferencia de ciberseguridad, escuché que los teléfonos no tienen una protección directa sobre los códigos QR, que simplemente los leen y abren lo que traen dentro. Me puse como misión encontrar una herramienta que hiciera con los QR lo que VirusTotal hace con los enlaces, y no encontré nada tan sólido.

La idea se quedó conmigo. Cada vez que veía un código QR en un restaurante, en un parqueo o en un cartel, pensaba en lo complicado que sería tomarle una foto, pasarla por un lector, extraer el enlace y recién entonces revisarlo. Y además, ¿qué pasaría si el enlace era malicioso pero todavía nadie lo había reportado? Ninguna lista lo iba a detectar.

Mi inspiración para la forma de la aplicación fue el lector de códigos QR que trae el iPhone en su centro de control, junto a los botones de sonido y de brillo. Lo uso mucho, y justamente por eso soy consciente de que no revisa nada. Me imaginé una herramienta igual de simple y de rápida de abrir, pero con la misión de leer los códigos de forma segura y de decirme a dónde me llevan antes de que yo llegue ahí.

Así nació QRisk: de una curiosidad, de un temor y de una necesidad personal de tener un escudo entre las nuevas tecnologías y nuestra seguridad. Este trabajo documenta el camino de convertir esa idea en un prototipo funcional.
